%%%%%%%%%%%%%%%%%%%%%%%%%%%%%%%%%%%%%%%%%%%%%%%%%%%%%%%%%%%%%%%%%%%%%%
% Lớp 10 | Thống kê | Các số đặc trưng đo mức độ phân tán của mẫu số liệu
% Dạng 0D6?4-4: Phương sai, độ lệch chuẩn của mẫu số liệu
% Số câu: 19  (sinh tự động bởi tach_cau.py; từ nay thầy sửa trực tiếp file này)
%%%%%%%%%%%%%%%%%%%%%%%%%%%%%%%%%%%%%%%%%%%%%%%%%%%%%%%%%%%%%%%%%%%%%%

%%%%%%%%%%%%%%%% TRẮC NGHIỆM 4 PHƯƠNG ÁN %%%%%%%%%%%%%%%%

% nguồn: MaHoaNH_TracNghiem_Lop10
\begin{ex}%[0D6N4-4]%[ma:6d8e8aad]
	Cho phương sai của các số liệu bằng $4$. Tìm độ lệch  chuẩn.
	\choice 
		{$4$} 
		{\True  $2$}
		{$16$} 
		{$8$} 
	\loigiai{ 
		Ta có độ lệch  chuẩn là căn bậc hai của phương sai. Suy ra \[s_x=\sqrt{s_x^2}=\sqrt{4}=2.\]
	} 
\end{ex}

% nguồn: MaHoaNH_TracNghiem_Lop10
\begin{ex}%[0D6N4-4]%[ma:1e82b07c]
	Theo kết quả thống kê điểm thi giữa kỳ $2$ môn toán khối $11$ của một trường THPT, người ta tính được phương sai của bảng thống kê đó là  $s_x^2=0{,}573$. Độ lệch chuẩn của bảng thống kê đó xấp xỉ bằng
	\choice 
		 {$0{,}812$}
		{\True  $0{,}757$}
		 {$0{,}936$}
		 {$0{,}657$}
	 \loigiai{
		Ta  có \[s_x=\sqrt{s_x^2}=\sqrt{0{,}573}\approx 0{,}757.\]
	} 
\end{ex}

% nguồn: MaHoaNH_TracNghiem_Lop10
\begin{ex}%[0D6N4-4]%[ma:3d7d7e62]
	Độ  lệch chuẩn là
	\choice 
		{\True căn  bậc hai của phương sai}
		{bình  phương của phương sai}
		{một nửa của  phương sai}
		{không phải các công  thức trên}
	\loigiai{ 
		Ta có độ lệch chuẩn là $s_x=\sqrt{s_x^2}$ nên độ lệch chuẩn là căn bậc hai  của phương sai.
	} 
\end{ex}

% nguồn: MaHoaNH_TracNghiem_Lop10
\begin{ex}%[0D6N4-4]%[ma:eb4909e6]
	Phương sai của  dãy số $2$; $3$; $4$; $5$; $6$ là
	\choice 
		{$s_x^2=4$} 
		 {$s_x^2=\sqrt{2}$}
		{\True  $s_x^2=2$}
		{$s_x^2=-2$} 
	\loigiai{ 
		 Ta có $\overline{x}=\dfrac{2+3+4+5+6}{5}=4$ nên suy ra \[s_x^2=\dfrac{1}{5}\left[(2-4)^2+(3-4)^2+(5-4)^2+(6-4)^2\right]=2.\]
	} 
\end{ex}

% nguồn: MaHoaNH_TracNghiem_Lop10
\begin{ex}%[0D6N4-4]%[ma:e36effc2]
	Cho phương sai của một mẫu số liệu là $0{,}6$ thì độ lệch chuẩn của mẫu  số liệu đó là
	 \choice
		{\True  $\dfrac{\sqrt{15}}{5}$}
		 {$\dfrac{\sqrt{3}}{5}$}
		{$\dfrac{1}{\sqrt{5}}$} 
		 {$\dfrac{3}{5}$}
	 \loigiai{
		Ta có $s^2=0{,}6=\dfrac{3}{5}$. Suy ra độ lệch chuẩn là $\sqrt{\dfrac{3}{5}}=\dfrac{\sqrt{15}}{5}$. 
	} 
\end{ex}

% nguồn: MaHoaNH_TracNghiem_Lop10
\begin{ex}%[0D6H4-4]%[ma:d9929c55]
	Bảng số liệu sau cho biết thời gian làm bài  tính bằng phút của $50$ học sinh.
	\begin{longtable}{|c|c|c|c|c|c|c|c|c|c|c|c|}
		\hline 
		Thời  gian&$3$&$4$&$5$&$6$&$7$&$8$&$9$&$10$&$11$&$12$&\\
		\hline 
		Số học sinh&$1$&$3$&$4$&$7$&$8$&$9$&$8$&$5$&$3$&$2$&$N=50$\\ 
		 \hline
	\end{longtable}
	Phương sai của mẫu  số liệu thống kê trên bằng
	\choice 
		{$7{,}68$} 
		{$2{,}13$} 
		{$4{,}63$} 
		{\True $4{,}54$} 
	\loigiai{ 
		Thời gian làm  bài trung bình của $50$ học sinh
		$$\overline{x}=\dfrac{1}{50}(3\cdot 1+4\cdot 3+5\cdot 4+6\cdot 7+7\cdot 8+8\cdot 9+9\cdot  8+10\cdot 5+11\cdot 3+12\cdot 2)= 7{,}68.$$
		Phương sai của mẫu số liệu \\$S^2=\dfrac{1\cdot 3^2+3\cdot4^2+4\cdot5^2+7\cdot6^2+8\cdot7^2+9\cdot8^2+8\cdot9^2+5\cdot10^2+3\cdot11^2+2\cdot12^2}{50}-7{,}68^2$\\$\approx 4{,}54$. 
	} 
\end{ex}

% nguồn: MaHoaNH_TracNghiem_Lop10
\begin{ex}%[0D6H4-4]%[ma:3a405b53]
	Theo dõi thời gian làm một bài toán (tính  bằng phút) của $40$ học sinh, giáo viên lập được bảng sau
	\begin{center}
		\begin{tabular}{|c|c|c|c|c|c|c|c|c|c|c|}
			\hline Thời gian $(x)$ & $4$ & $5$ & $6$ & $7$ & $8$ & $9$ & $10$ & $11$  & $12$ & \\
			\hline  Tần số $(n)$ & $6$ & $3$ & $4$ & $2$ & $7$ & $5$ & $5$ & $7$ & $1$ & $N=40$\\
			 \hline
		\end{tabular}
	\end{center}
	Phương sai của mẫu số liệu trên gần với số nào nhất? 
	\choice 
		 {\True $6$}
		 {$12$}
		 {$40$}
		{$9$} 
	 \loigiai{
		Ta có giá  trị trung bình của mẫu số liệu là
		 \[\overline{x} = \dfrac{1}{N}\left( n_1x_1+n_2x_2+ \cdots + n_kx_k\right) =\dfrac{317}{40}.\]
		Suy ra phương sai  của mẫu số liệu là
		\[s^2 = \dfrac{1}{N}\left[n_1\left(x_1-\overline{x}\right)^2 + n_2\left(x_2-\overline{x}\right)^2 + \cdots + n_k\left(x_k-\overline{x}\right)^2\right] \approx  6.\]
	} 
\end{ex}

% nguồn: MaHoaNH_TracNghiem_Lop10
\begin{ex}%[0D6H4-4]%[ma:e561c6b7]
	Điểm thi  cuối học kì $2$ gồm $6$ môn học của một học sinh được cho vào bảng sau
	\begin{center}
		\begin{tabular}{|c|c|c|c|c|c|}
			\hline Toán & Lí & Hóa &  Sinh & Văn & Anh \\
			\hline 10 & 9 & 8,5 & 8 & 7 &  5 \\
			\hline 
		\end{tabular}
	\end{center}
	Độ lệch chuẩn  $s$ của mẫu số liệu trên gần nhất với giá trị nào sau đây?
	\choice 
		{\True  $1{,}59$}
		{$2{,}53$} 
		{$3{,}04$} 
		{$1{,}74$} 
	\loigiai{ 
		Trung bình  cộng của mẫu số liệu
		 $$\overline{x}=\dfrac{10+9+8{,}5+8+7+5}{6}=\dfrac{95}{12}.$$
		Phương sai của  mẫu số liệu
		 $$s^2=\dfrac{\left(10-\dfrac{95}{12}\right)^2+\left(9-\dfrac{95}{12}\right)^2+\left(8{,}5-\dfrac{95}{12}\right)^2+\left(8-\dfrac{95}{12}\right)^2+\left(7-\dfrac{95}{12}\right)^2+\left(5-\dfrac{95}{12}\right)^2}{6}=\dfrac{365}{144}.$$
		Độ lệch  chuẩn của mẫu số liệu
		 $$s=\sqrt{s^2}=\sqrt{\dfrac{365}{144}}\approx 1{,}59.$$
	 }
\end{ex}

% nguồn: MaHoaNH_TracNghiem_Lop10
\begin{ex}%[0D6H4-4]%[ma:8f72a5a0]
	Điểm thi của lớp $10$C của một trường Trung học Phổ Thông được trình bày ở bảng phân bố tần  số sau
	\begin{center}
		\begin{tabular}{|c|c|c|c|c|c|c|c|}
			\hline Điểm thi & $5$ & $6$ & $7$  & $8$ & $9$ & $10$ & \\
			\hline Tần số & $7$ & $5$  & $10$ & $12$ & $4$ & $2$ & $n=40$\\
			\hline 
		\end{tabular}
	\end{center}
	Phương sai của bảng phân  bố tần số đã cho là
	 \choice
		 {$0{,}94$}
		 {$3{,}94$}
		{$2{,}94$} 
		{\True  $1{,}94$}
	 \loigiai{
		Trong  dãy số liệu về điểm thi của lớp $10$C ta có
		 \[\overline{x} = \dfrac{1}{n}\left( n_1x_1+n_2x_2+ \cdots + n_6x_6\right) =\dfrac{1}{40}\left(7\cdot 5+5\cdot 6+10\cdot 7+12\cdot 8+4\cdot 9+2\cdot 10\right)=7{,}175.\]
		Suy ra phương  sai là
		\begin{align*}
			s^2 &= \dfrac{1}{n}\left[n_1\left(x_1-\overline{x}\right)^2 + n_2\left(x_2-\overline{x}\right)^2 + \cdots + n_6\left(x_6-\overline{x}\right)^2\right] 
			\\  &= \dfrac{1}{40}\left[7\left(5-7{,}175\right)^2 + 5\left(6-7{,}175\right)^2 + 10\left(7-7{,}175\right)^2 + 12\left(8-7{,}175\right)^2 \right.
			\\ & \quad \left.  + 4\left(9-7{,}175\right)^2 + 2\left(10-7{,}175\right)^2 \right]
			 \\ & \approx 1{,}94.
		\end{align*}
	 }
\end{ex}


%%%%%%%%%%%%%%%% ĐÚNG/SAI %%%%%%%%%%%%%%%%

% nguồn: MaHoaNH_DungSai_Lop10
\begin{ex}%[0D6H4-4]%[ma:dfac141f]
	Mẫu số liệu dưới đây thống kê thời gian chờ xe bus (đơn vi: phút) của $10$ học sinh ở cùng một bến: $1, 4, 5, 6, 6, 8, 10, 11,  12, 25$. Khi đó
	\choiceTF 
		{\True Thời gian  chờ xe bus trung bình của $10$ học sinh trên là $8{,}8$ phút}
		 {Mốt của mẫu số liệu trên bằng $25$}
		{\True Giá trị ngoại lệ  của mẫu số liệu trên là $25$}
		{\True Độ lệch chuẩn về thời gian chờ xe bus của $10$ học sinh trên là $6{,}27$  phút}
	 \loigiai{
		\begin{itemchoice}
			\itemch{\bf  Đúng.}\\
			Thời gian chờ xe bus trung bình của  $10$ học sinh trên là \[\overline{x}=\dfrac{1+4+5+6+6+8+10+11+12+25}{10}=8{,}8\,(\text{phút}).\]
			\itemch{\bf  Sai.}\\
			Dựa vào mẫu số liệu ta thấy thời gian chờ xe bus $6$ phút có $2$ học sinh, những  thời gian còn lại chỉ có $1$ học sinh.\\
			Do đó mốt của mẫu số  liệu là $M_0=6$.
			 \itemch{\bf Đúng.}\\
			Sắp xếp  mẫu số liệu theo thứ tự không giảm $1, 4, 5, 6, 6, 8, 10, 11, 12, 25$.\\
			Cỡ  mẫu $n=10$.\\
			Tứ phân vị thứ  nhất là $Q_1=x_3=5$.\\
			Tứ  phân vị thứ ba là $Q_3=x_8=11$.\\
			Khoảng tứ phân vị $\Delta  Q=Q_3-Q_1=11-5=6$.\\
			Ta có $\Delta Q-1{,}5\cdot Q_1=6-1{,}5\cdot 5=-1{,}5$ và $\Delta Q+1{,}5\cdot Q_3=6+1{,}5\cdot 11=22{,}5$.\\ 
			Ta thấy  $25>22{,}5$ nên $25$ là giá trị ngoại lệ.
			\itemch{\bf  Đúng.}\\
			 Phương sai $S^2=\dfrac{1^2+4^2+5^2+6^2+6^2+8^2+10^2+11^2+12^2+25^2}{10}-8{,}8^2=39{,}36$.\\
			Độ  lệch chuẩn $S=\sqrt{S^2}=\sqrt{39{,}36}\approx 6{,27}$.
		\end{itemchoice}
	 }
\end{ex}

% nguồn: MaHoaNH_DungSai_Lop10
\begin{ex}%[0D6H4-4]%[ma:bc9a41f7]
	 Điều tra thu nhập của $40$ hộ gia đình ở một bản được số liệu như sau
	\begin{center}
		\begin{tabular}{|c|c|c|c|c|c|c|c|c|c|}
			\hline Thu nhập (triệu đồng/năm) & $4$ & $4{,}5$ & $5$ & $5{,}5$ & $6$ & $6{,}5$ & $7$ &  $7{,}5$ & $8$ \\
			\hline Số hộ gia đình & $1$ & $3$ & $4$ & $6$ &  $8$ & $7$ & $6$ & $3$ & $2$ \\
			\hline 
		\end{tabular}
	\end{center}
	Các mệnh đề  sau đúng hay sai?
	\choiceTF 
		{\True Cỡ mẫu là  $40$}
		{Mốt của mẫu số liệu  là $8$}
		{\True  Số trung bình của mẫu số liệu trên là $6{,}1125$}
		{Phương sai  của mẫu số liệu trên là $1{,}9561$}
	 \loigiai{
		\begin{itemchoice}
			\itemch Đúng. Từ giả thiết ta có  cỡ mẫu là $40$.
			\itemch Sai. Có $8$ hộ gia đình thu nhập  $6$ triệu đồng/năm nên mốt là $6$.
			\itemch Đúng. Số trung bình của mẫu số liệu trên là  \[\overline{x}=\dfrac{4\cdot 1+4{,}5\cdot 3+5\cdot 4+5{,}5\cdot 6+6\cdot 8+6{,}5\cdot 7+7\cdot 6+7{,}5\cdot 3 +8\cdot 2}{40}=6{,}1125.\]
			\itemch Sai. Phương sai của mẫu số liệu trên  là
			\begin{align*}
				s^2 &= \dfrac{\left(4^2\cdot 1+4{,}5^2\cdot 3+5^2\cdot 4+5{,}5^2\cdot 6+6^2\cdot 8+6{,}5^2\cdot  7+7^2\cdot 6+7{,}5^2\cdot 3 +8^2\cdot 2\right)}{40}-6{,}1125^2
				\\ &\approx 0{,}9561. 
			\end{align*}
		\end{itemchoice}
	} 
\end{ex}

% nguồn: MaHoaNH_DungSai_Lop10
\begin{ex}%[0D6H4-4]%[0D6H3-5]%[0D6H4-2]%[0D6H4-4]%[ma:fd3e0f8f]
	Có $ 100 $ học sinh lớp 10 tham dự kì thi học sinh giỏi Toán cấp cụm (thang điểm $ 20 $). Kết  quả cho trong bảng sau
	\begin{longtable}{|c|c|c|c|c|c|c|c|c|c|c|c|}
		\hline Điểm $ (x) $ & $9$ & $10$ & $11$ & $12$ & $13$ & $14$ & $15$ & $16$ & $17$ & $18$& $19$\\ 
		\hline Tần số & $3$ & $4$ & $5$ & $3$ &  $7$ & $15$ & $20$ & $14$ & $19$ & $6$& $4$\\
		\hline 
	\end{longtable}
	 \choiceTF
		{\True Điểm số  trung bình của của $100$ học sinh là $ 14{,}9 $}
		 {Mốt của mẫu số liệu là $20$}
		{Khoảng biến thiên của mẫu  số liệu là $1$}
		{\True Độ lệch chuẩn của mẫu số liệu là $  2{,}38 $}
	\loigiai{ 
		\begin{itemize}
			\item  Điểm số trung bình là $$\overline{x}=\dfrac{9\cdot3+10\cdot4+11\cdot5+12\cdot3+13\cdot7+14\cdot15+15\cdot20+16\cdot14+17\cdot19+18\cdot6+19\cdot4}{100}=14{,}9.$$
			\item Mốt của mẫu số liệu là  $15$ vì $15$ điểm có tần số lớn nhất bằng $20$.
			\item Khoảng biến thiên của mẫu số liệu là $19-9=10$. 
			\item  Ta có $ 3(9-14{,}9)^2+4(10-14{,}9)^2+5(11-14{,}9)^2+3(12-14{,}9)^2+7(13-14{,}9)^2+15(14-14{,}9)^2+20(15-14{,}9)^2+14(16-14{,}9)^2+19(17-14{,}9)^2+6(18-14{,}9)^2+4(19-14{,}9)^2=565 $.\\
			Phương sai của mẫu số  liệu là $ s^2=\dfrac{565}{100}=5{,}65 $.\\
			Suy ra độ  lệch chuẩn của mẫu số liệu là $ s=\sqrt{s^2}=\sqrt{5{,}65}\approx 2{,}38 $.
		\end{itemize}
		\begin{itemchoice}
			\itemch Đúng. 
			\itemch  Sai.
			\itemch Sai. 
			\itemch Đúng. 
		\end{itemchoice}
	} 
\end{ex}


%%%%%%%%%%%%%%%% TRẢ LỜI NGẮN %%%%%%%%%%%%%%%%

% nguồn: MaHoaNH_TraLoiNgan_Lop10
\begin{bt}%[0D6H4-4]%[ma:3c29b7d0]
	Cho mẫu số liệu thống kê $\{1;2;3;4;5;6;7;8;9\}$. Tính độ  lệch chuẩn của mẫu số liệu trên (làm tròn đến hàng phần trăm).
	 \shortans{2,58}
	 \loigiai{
		Ta có giá trị trung bình của  mẫu số liệu là \[\overline{x}=\dfrac{1+2+3+4+5+6+7+8+9}{9}=5.\]
		Suy ra độ lệch  chuẩn là \[s=\sqrt{\dfrac{1}{9}\left(1^2+2^2+3^2+4^2+5^2+6^2+7^2+8^2+9^2\right)-5^2}\approx 2{,}58.\]
	} 
\end{bt}

% nguồn: MaHoaNH_TraLoiNgan_Lop10
\begin{bt}%[0D6H4-4]%[ma:68facdc4]
	Để biết cây đậu phát triển như thế nào sau khi gieo hạt, bạn Châu gieo $5$ hạt đậu vào $5$ chậu riêng biệt và cung cấp cho chúng lượng nước, ánh sáng như nhau. Sau  $2$ tuần, $5$ hạt đậu đã nảy mầm và phát triển thành $5$ cây con. Bạn Châu đo chiều cao từ rễ đến ngọn của mỗi cây (đơn vị mm) và ghi kết quả là mẫu số liệu sau
	 \[112 \quad 102 \quad 106 \quad 94 \quad 101 \]
	 Tìm độ lệch chuẩn của mẫu số liệu trên (kết quả làm tròn đến hàng phần trăm).
	 \shortans{$5{,}93$}
	\loigiai{ 
		Từ bảng số liệu ta có số  trung bình là \[\overline{x}=\dfrac{112+102+106+94+101}{5}=103.\]
		Suy ra ta tính được phương sai và độ lệch chuẩn lần  lượt là \[s^2=35{,}2 \text{ và } s\approx 5{,}93.\]
	} 
\end{bt}

% nguồn: MaHoaNH_TraLoiNgan_Lop10
\begin{bt}%[0D6H4-4]%[ma:6f0a3d6c]
	Số liệu thống kê tỉ lệ $(\%)$ tốt nghiệp THPT của một địa phương từ năm học $2001 - 2002$ đến năm học $2016 - 2017$ được cho như  sau
	\begin{center}
		\begin{tabular}{|c|c|c|c|c|c|c|c|}
			\hline  $98{,}82$ & $97{,}46$ & $99{,}19$ & $98{,}90$ & $98{,}65$ & $79{,}51$ & $85{,}06$ & $86{,}18$ \\
			\hline $98{,}68$ & $99{,}23$  & $99{,}93$ & $99{,}34$ & $99{,}74$ & $93{,}08$ & $97{,}34$ & $97{,}82$ \\
			 \hline
		\end{tabular}
	\end{center}
	Tính phương sai của mẫu số liệu. 
	 \shortans{$37{,}0$}
	 \loigiai{
		Phương sai là \[s^2=\dfrac{1}{16}\left[\left(x_1-\overline{x}\right)^2+  \left(x_2-\overline{x}\right)^2+\cdots + \left(x_{16}-\overline{x}\right)^2\right]\approx 37{,}03.\]
	} 
\end{bt}

% nguồn: MaHoaNH_TraLoiNgan_Lop10
\begin{bt}%[0D6H4-4]%[ma:673c7ff6]
	Sản lượng lúa (đơn vị là tạ) của $40$ thửa ruộng thí nghiệm có cùng diện tích được trình bày trong bảng số liệu  sau
	\begin{center}
		\begin{tabular}{|c|c|c|c|c|c|c|c|c|c|}
			\hline Sản lượng & $20$ & $21$ & $22$ &  $23$ & $24$ & \\
			\hline Tần số & $5$ & $8$ & $11$ & $10$ & $6$ & $N=40$ \\ 
			\hline 
		\end{tabular}
	\end{center}
	Tìm phương sai của mẫu số liệu trên (làm  tròn đến hàng phần trăm).
	 \shortans{1,54}
	 \loigiai{
		Ta có giá trị trung bình của mẫu số liệu là \[\overline{x}=\dfrac{20\cdot 5+21\cdot 8+22\cdot  11+23\cdot 10+24\cdot 6}{40}=22{,}1.\]
		Suy ra phương sai là \[s^2=\dfrac{1}{40}\left(20^2\cdot 5+21^2\cdot 8+22^2\cdot 11+23^2\cdot  10+24^2\cdot 6\right)-22{,}1^2\approx 1{,}54.\]
	} 
\end{bt}

% nguồn: MaHoaNH_TraLoiNgan_Lop10
\begin{bt}%[0D6H4-4]%[ma:a9a891e6]
	Bảng  số liệu sau thống kê nhiệt độ tại Thành phố Hồ Chí Minh trong một lần đo vào một ngày của năm $2021$ như sau
	\begin{center}
		\begin{tabular}{|c|c|c|c|c|c|c|c|c|c|c|c}
			\hline  Giờ đo & $1$h & $4$h & $7$h & $10$h & $13$h & $16$h & $19$h & $22$h \\
			\hline Nhiệt độ ($^\circ \mathrm{C}$) & $27$ & $26$  & $28$ & $32$ & $34$ & $35$ & $30$ & $28$ \\
			 \hline
		\end{tabular}
	\end{center}
	Tìm độ lệch chuẩn của  mẫu số liệu đã cho (làm tròn kết quả đến hàng phần trăm).
	\shortans{$3{,}12$} 
	\loigiai{ 
		Ta có số trung bình là  \[\overline{x}=\dfrac{27+26+28+32+34+35+30+28}{8}=30^\circ C.\]
		Suy ra độ lệch chuẩn là \[s=\sqrt{\dfrac{1}{8}\left[(27-30)^2+(26-30)^2+(28-30)^2+\cdots+(28-30)^2\right] }\approx 3{,}12.\] 
	} 
\end{bt}

% nguồn: MaHoaNH_TraLoiNgan_Lop10
\begin{bt}%[0D6H4-4]%[ma:6a70cb10]
	Trong $5$ lần nhảy xa, hai bạn Hùng và Trung có kết quả (đơn vị mét) lần lượt  là
	\begin{center}
		\begin{tabular}{|c|c|c|c|c|c|c|c|c|}
			\hline Hùng & $2{,}4$ & $2{,}6$ & $2{,}4$ & $2{,}5$ &  $2{,}6$ \\
			\hline Trung &  $2{,}4$ & $2{,}5$ & $2{,}5$ & $2{,}5$ & $2{,}6$ \\
			\hline 
		\end{tabular}
	\end{center}
	 Hỏi phương sai mẫu số liệu của Hùng gấp của Trung bao nhiêu lần?
	 \shortans{2}
	\loigiai{ 
		Phương sai của mẫu số liệu là 
		\begin{itemize}
			\item Hùng $s^2=\dfrac{1}{5}\left(2{,}4^2+2{,}6^2+2{,}4^2+2{,}5^2+2{,}6^2\right)-2{,}5^2\approx 0{,}008$. 
			 \item Trung $s^2=\dfrac{1}{5}\left(2{,}4^2+2{,}5^2+2{,}5^2+2{,}5^2+2{,}6^2\right)-2{,}5^2\approx 0{,}004$.
		\end{itemize}
		Vậy phương sai mẫu số liệu của Hùng gấp của Trung $2$ lần. 
	} 
\end{bt}

% nguồn: MaHoaNH_TraLoiNgan_Lop10
\begin{bt}%[0D6V4-4]%[ma:4822a023]
	Bài thi Tiếng Anh gồm có $100$ câu trắc nghiệm, mỗi đáp án chọn đúng được $1$ điểm,  chọn sai $0$ điểm. Kết quả kiểm tra của lớp 10A được thống kê như sau
	\begin{longtable}{|c|c|c|c|c|c|c|c|c|c|c|c|c|c|c|}
		\hline 
		 54&67&87&23&54&76&15&64&74&35&65&60&62&50&46\\
		 \hline
		 58&61&49&49&58&59&59&79&82&100&95&64&55&38&72\\
		\hline 
	\end{longtable}\noindent
	Gọi $Q_1$, $Q_3$, $S$ lần lượt là tứ phân vị thứ  nhất, tứ phân vị thứ ba và phương sai của mẫu số liệu trên. Tính giá trị của biểu thức $P=Q_1+Q_3+S$ (làm tròn kết quả đến hàng đơn vị).
	\shortans{583} 
	 \loigiai{
		Mẫu số liệu được sắp xếp theo thứ tự không  giảm
		\begin{longtable}{|c|c|c|c|c|c|c|c|c|c|c|c|c|c|c|}
			\hline 
			$15$&$23$&$35$& $38$&$46$&$49$&$49$&$50$&$54$&$54$&$55$& $58$& $58$& $59$&  $59$\\
			\hline 
			$60$&$61$&$62$&$64$&$64$&$65$&$67$&$72$&$74$&$76$&$79$&$82$&$87$& $95$&  $100$\\
			 \hline
		\end{longtable}\noindent
		Cỡ mẫu $n=30$.\\ 
		Tứ phân  vị thứ nhất $Q_1=x_8=50$.\\
		Tứ  phân vị thứ ba $Q_3=x_{23}=72$.\\
		Điểm  trung bình $\overline{x}=\dfrac{15+23+35+\ldots+87+95+100}{30}=\dfrac{178}{3}$.\\
		 Phương sai $S=\dfrac{15^2+23^2+35^2+\ldots+87^2+95^2+100^2}{30}-\left(\dfrac{178}{3} \right)^2=\dfrac{20\,761}{45}$.\\
		Vậy  $P=Q_1+Q_3+S=50+72+\dfrac{20\,761}{45}\approx 583$.
	} 
\end{bt}
